% ─────────────────────────────────────────────────────────────────────────────
%
% raspi_docker_config.tex
%
% Documentación de la configuración de despliegue en Raspberry Pi:
%   - Entorno Docker: imagen única para PC (amd64) y Raspberry Pi (arm64)
%   - Diferencias de configuración entre ambas plataformas
%   - Servicio systemd nav-mapper.service para arranque automático
%   - Procedimiento de instalación y operación del servicio
%
% ─────────────────────────────────────────────────────────────────────────────

\section{Configuración de despliegue en Raspberry Pi}
\label{sec:raspi-deploy}

El sistema se desarrolla y prueba en una PC de escritorio con arquitectura
\texttt{amd64} y se despliega en una Raspberry Pi con arquitectura
\texttt{arm64}. Para garantizar paridad de entorno entre ambas plataformas
se utiliza Docker como capa de aislamiento: una única imagen base se
construye de forma nativa en cada arquitectura, y las diferencias de
plataforma se resuelven íntegramente dentro del \texttt{Dockerfile} en tiempo
de construcción, sin afectar al código fuente de los paquetes ROS2.

% ─────────────────────────────────────────────────────────────────────────────
\subsection{Entorno Docker}
\label{sec:docker-env}
% ─────────────────────────────────────────────────────────────────────────────

\subsubsection{Imagen base y construcción multi-plataforma}

Ambas plataformas parten de la imagen oficial \texttt{ros:humble-ros-base},
que provee ROS~2 Humble sobre Ubuntu~22.04. La imagen se construye de forma
nativa en cada máquina (no se usa emulación QEMU), por lo que el argumento
\texttt{TARGETARCH} es rellenado automáticamente por el daemon de Docker con
el valor correcto: \texttt{amd64} en la PC y \texttt{arm64} en la Raspberry
Pi.

La imagen resultante en la Raspberry Pi se etiqueta como
\texttt{tesis-ros2:humble-cpu-arm64}; en la PC como
\texttt{tesis-ros2:humble-cpu}.

\subsubsection{Capas comunes a ambas plataformas}

Las siguientes dependencias se instalan en todas las arquitecturas:

\begin{itemize}
  \item Herramientas de construcción: \texttt{colcon}, \texttt{vcstool},
        \texttt{cmake}, \texttt{clang-format}.
  \item Paquetes ROS2: \texttt{realsense2-camera},
        \texttt{vision\_opencv}, \texttt{rmw\_cyclonedds\_cpp},
        \texttt{imu\_filter\_madgwick}, \texttt{rtabmap-ros},
        \texttt{rviz2}.
  \item Librería de acceso GPIO: \texttt{libgpiod-dev} (usada por el
        paquete \texttt{gpio\_rosbag\_controller} para disparar grabaciones
        mediante un pulsador físico).
\end{itemize}

\subsubsection{Diferencias exclusivas de la PC (amd64)}

En la PC se instalan herramientas gráficas que no tienen sentido en la
Raspberry Pi, donde no hay display conectado:

\begin{itemize}
  \item \texttt{ros-humble-rqt} y \texttt{ros-humble-rqt-image-view}:
        paneles de introspección de tópicos en tiempo real.
  \item Soporte \texttt{GTK3} / \texttt{Mesa GL}: necesario para las
        ventanas de OpenCV (\texttt{imshow}) y RViz2.
\end{itemize}

\subsubsection{Diferencias exclusivas de la Raspberry Pi (arm64)}

El paquete oficial de apt \texttt{ros-humble-realsense2-camera} enlaza
contra \texttt{librealsense2} versión~2.56, que internamente usa el
\emph{backend de kernel} (\texttt{V4L2/UVC}). Este backend requiere
permisos de root y en Raspberry Pi OS presenta problemas de inicialización
con la cámara D435i. Para resolverlo se compila \texttt{librealsense}
v2.57.6 desde su código fuente activando el \emph{RSUSB backend}, que
opera en espacio de usuario sobre \texttt{libusb-1.0}:

\begin{verbatim}
cmake .. \
  -DFORCE_RSUSB_BACKEND=ON \
  -DBUILD_EXAMPLES=OFF \
  -DBUILD_WITH_OPENMP=OFF
make -j2 && make install
\end{verbatim}

La biblioteca resultante queda en \texttt{/usr/local/lib/librealsense2.so.2.57.*}.
Para que el wrapper de ROS2 la utilice en lugar de la versión de apt se crea
un enlace simbólico sobre el archivo que instala el paquete:

\begin{verbatim}
ln -sf /usr/local/lib/librealsense2.so.2.57.* \
       /opt/ros/humble/lib/aarch64-linux-gnu/librealsense2.so.2.56
\end{verbatim}

Además se configura el orden de búsqueda de la carga dinámica para que
\texttt{/usr/local/lib} y \texttt{/usr/local/bin} tengan precedencia:

\begin{verbatim}
ENV PATH="/usr/local/bin:${PATH}"
ENV LD_LIBRARY_PATH="/usr/local/lib:${LD_LIBRARY_PATH}"
\end{verbatim}

Esta misma configuración se replica en \texttt{/etc/bash.bashrc} para que
también esté disponible en sesiones interactivas.

Por limitaciones de la Pi, la compilación de librealsense usa sólo
\texttt{-j2} (dos hilos) para no agotar la memoria RAM durante el build de
la imagen.

\subsubsection{Usuario en el contenedor}

El contenedor crea internamente un usuario cuyo nombre, UID y GID se
inyectan en tiempo de construcción mediante argumentos \texttt{ARG}:

\begin{verbatim}
docker compose build
  --build-arg USERNAME=tomasdea
  --build-arg USER_UID=$(id -u)
  --build-arg USER_GID=$(id -g)
\end{verbatim}

En la práctica estos valores son leídos desde las variables de entorno del
\texttt{docker-compose.raspi.yml}. Esto asegura que los archivos creados
dentro del contenedor (bags, logs, archivos de build) tengan el mismo
propietario que el usuario del host, evitando problemas de permisos al
acceder a los volúmenes montados.

\subsubsection{Volúmenes y dispositivos montados}

El workspace del host se monta en el mismo path dentro del contenedor:

\begin{verbatim}
${HOST_WS}:${CONTAINER_WS}
\end{verbatim}

De este modo \texttt{HOST\_WS} y \texttt{CONTAINER\_WS} apuntan al mismo
directorio físico (\texttt{/home/tomasdea/Tesis/develop/Fiuba-Tesis} en la
Pi), lo que simplifica la depuración: editar un archivo en el host es
equivalente a hacerlo en el contenedor.

En la Raspberry Pi también se montan los dispositivos necesarios para el
hardware:

\begin{table}[H]
\centering
\begin{tabular}{ll}
\toprule
\textbf{Recurso montado} & \textbf{Finalidad} \\
\midrule
\texttt{/dev} (completo) & Acceso a la cámara RealSense y al UART \\
\texttt{/dev/bus/usb}    & Acceso USB explícito (RSUSB backend) \\
\texttt{/sys}            & Información de topología de hardware \\
\texttt{/run/udev}       & Eventos de conexión de dispositivos \\
\texttt{/dev/gpiochip0}  & GPIO para el pulsador de grabación \\
\bottomrule
\end{tabular}
\caption{Dispositivos del host montados en el contenedor (Raspberry Pi).}
\label{tab:raspi-mounts}
\end{table}

En la PC de desarrollo no se monta \texttt{/dev/gpiochip0} ni se agrega el
grupo \texttt{gpio} al contenedor, ya que esa funcionalidad es exclusiva de
la Pi. En cambio, sí se monta \texttt{/tmp/.X11-unix} para permitir la
apertura de ventanas gráficas (RViz2, \texttt{imshow}) en el display del
host.

\subsubsection{Contexto de shell y variable WS\_PATH}

El archivo \texttt{docker/context.sh} se monta como
\texttt{/etc/profile.d/context.sh} en el contenedor (solo lectura).
\texttt{/etc/profile.d/} es un directorio especial de Linux que todos los
shells de \emph{login} (\texttt{bash -l}) ejecutan automáticamente al
iniciar. El archivo realiza dos tareas:

\begin{enumerate}
  \item Detecta si el proceso está corriendo dentro de un contenedor Docker
        (consulta \texttt{/.dockerenv}, \texttt{/proc/1/cgroup} y la
        variable \texttt{RUNNING\_IN\_CONTAINER}) y exporta
        \texttt{SHELL\_CONTEXT=container} o \texttt{host} según corresponda.
  \item Sourcea \texttt{\$WS\_PATH/tools/dev/context\_env.sh}, que define
        aliases y helpers de desarrollo (\texttt{tesis-run}, \texttt{tesis-build},
        etc.) y personaliza el prompt del bash con el tag \texttt{[docker]}
        o \texttt{[local]}.
\end{enumerate}

La variable \texttt{WS\_PATH} se inyecta desde el \texttt{docker-compose}
en tiempo de ejecución mediante:

\begin{verbatim}
environment:
  - WS_PATH=${CONTAINER_WS}
\end{verbatim}

y también a través de \texttt{BASH\_ENV=/etc/profile.d/context.sh}, que
hace que el archivo de contexto se sourcea incluso en shells
\emph{no-interactivos} (importante para el servicio systemd).

% ─────────────────────────────────────────────────────────────────────────────
\subsection{Servicio systemd \texttt{nav-mapper}}
\label{sec:systemd-service}
% ─────────────────────────────────────────────────────────────────────────────

\subsubsection{¿Qué es systemd?}

\textbf{systemd} es el gestor de sistema e inicialización (\emph{init system})
estándar en las distribuciones Linux modernas, incluyendo Raspberry Pi OS.
Es el primer proceso que arranca el kernel (PID~1) y es responsable de
iniciar todos los demás procesos del sistema en el orden correcto.

Un \textbf{servicio} (o \emph{unit} de tipo \texttt{service}) es un
archivo de texto con extensión \texttt{.service} que describe cómo
iniciar, detener, reiniciar y monitorear un proceso. systemd gestiona el
ciclo de vida completo del proceso: lo arranca al encender el equipo, lo
reinicia si falla, captura su salida estándar y la envía al journal del
sistema (consultable con \texttt{journalctl}).

La alternativa sin systemd sería entrar por SSH a la Pi cada vez que
se enciende, navegar al directorio, sourcear ROS2 y ejecutar el launch
manualmente. El servicio elimina ese paso y hace el sistema autónomo.

\subsubsection{Estructura de un archivo \texttt{.service}}

Un archivo \texttt{.service} tiene tres secciones obligatorias:

\begin{description}
  \item[\texttt{[Unit]}] Metadatos y dependencias. Define qué otros
    servicios deben estar activos antes de que este pueda arrancar
    (\texttt{After=}, \texttt{Requires=}).
  \item[\texttt{[Service]}] Configuración del proceso: usuario que lo ejecuta,
    variables de entorno, comando de arranque (\texttt{ExecStart}),
    política de reinicio.
  \item[\texttt{[Install]}] A qué \emph{target} (equivalente a un runlevel)
    pertenece el servicio. \texttt{multi-user.target} equivale al modo
    de operación normal con red pero sin entorno gráfico.
\end{description}

\subsubsection{Archivo \texttt{nav-mapper.service}}

El archivo de servicio del proyecto se encuentra en
\texttt{src/nav\_bringup/scripts/nav-mapper.service}. A continuación se
describe cada directiva:

\paragraph{Sección \texttt{[Unit]}}

\begin{verbatim}
After=network.target docker.service
Requires=docker.service
\end{verbatim}

\texttt{After} establece el \emph{orden} de arranque: systemd espera a que
\texttt{network.target} (red disponible, necesaria para DDS/ROS2) y
\texttt{docker.service} (daemon de Docker) estén activos antes de iniciar
este servicio. \texttt{Requires} declara una \emph{dependencia dura}: si
Docker no está disponible, el servicio no arranca.

\paragraph{Sección \texttt{[Service]}}

\begin{verbatim}
User=tomasdea
Group=tomasdea
\end{verbatim}

El proceso se ejecuta con el usuario \texttt{tomasdea} (no como root).
Esto es importante porque el contenedor Docker hereda el mismo usuario
y porque los archivos del workspace (bags, logs) quedarán con el
propietario correcto.

\begin{verbatim}
Environment="HOME=/home/tomasdea"
Environment="ROS_DOMAIN_ID=0"
Environment="WS_PATH=/home/tomasdea/Tesis/develop/Fiuba-Tesis"
\end{verbatim}

systemd no sourcea \texttt{.bashrc} ni \texttt{.profile}: el proceso
arranca con un entorno mínimo. Por eso todas las variables necesarias se
declaran explícitamente con \texttt{Environment=}. \texttt{WS\_PATH} es
consumida por \texttt{docker compose} para montar el volumen correcto y
también es inyectada dentro del contenedor.

\begin{verbatim}
ExecStart=/usr/bin/docker compose \
  -f ${WS_PATH}/docker-compose.raspi.yml \
  run --rm dev \
  bash -lc 'exec ros2 launch nav_bringup nav.launch.py \
            use_hw:=false use_gpio_recorder:=false'
\end{verbatim}

En lugar de sourcear ROS2 directamente en el host, el servicio levanta un
contenedor efímero usando el perfil de Raspberry Pi. El flag \texttt{--rm}
elimina el contenedor al terminar, evitando acumulación. El flag \texttt{-l}
de \texttt{bash} activa el modo \emph{login shell}, que ejecuta
\texttt{/etc/profile.d/context.sh} (montado por el compose), el cual
sourcea ROS2 y el workspace automáticamente. Esto es equivalente al flujo
manual de desarrollo con \texttt{tesis-run}.

\begin{verbatim}
Restart=on-failure
RestartSec=5s
TimeoutStopSec=30s
\end{verbatim}

Si el proceso termina con código de error (p.\,ej.\ la cámara no se detecta
al arrancar), systemd lo reinicia automáticamente tras 5~segundos.
\texttt{TimeoutStopSec=30s} da tiempo a que \texttt{ros2 bag record} (si
está activo) haga \texttt{flush} y cierre el bag limpiamente antes de que
systemd envíe \texttt{SIGKILL}.

\begin{verbatim}
StandardOutput=journal
StandardError=journal
SyslogIdentifier=nav-mapper
\end{verbatim}

Toda la salida del proceso (incluyendo los \texttt{RCLCPP\_INFO} de los
nodos ROS2) se envía al journal de systemd con el identificador
\texttt{nav-mapper}, lo que permite filtrarla fácilmente:
\texttt{journalctl -u nav-mapper -f}.

\subsubsection{Instalación del servicio en la Raspberry Pi}

Los pasos para instalar y habilitar el servicio son los siguientes:

\begin{enumerate}
  \item Copiar el archivo al directorio de servicios del sistema:
\begin{verbatim}
sudo cp src/nav_bringup/scripts/nav-mapper.service \
        /etc/systemd/system/
\end{verbatim}

  \item Recargar el índice de units de systemd para que detecte el nuevo
        archivo:
\begin{verbatim}
sudo systemctl daemon-reload
\end{verbatim}

  \item Habilitar el servicio para que arranque automáticamente en cada
        inicio del sistema:
\begin{verbatim}
sudo systemctl enable nav-mapper.service
\end{verbatim}

  \item Iniciarlo manualmente por primera vez (sin necesidad de reiniciar):
\begin{verbatim}
sudo systemctl start nav-mapper.service
\end{verbatim}

  \item Verificar que está corriendo:
\begin{verbatim}
systemctl status nav-mapper.service
\end{verbatim}

  \item Monitorear la salida en tiempo real:
\begin{verbatim}
journalctl -u nav-mapper -f
\end{verbatim}
\end{enumerate}

Para detener el servicio sin deshabilitarlo (útil durante desarrollo):
\begin{verbatim}
sudo systemctl stop nav-mapper.service
\end{verbatim}

Para deshabilitar el arranque automático permanentemente:
\begin{verbatim}
sudo systemctl disable nav-mapper.service
\end{verbatim}

\subsubsection{Permisos de Docker sin sudo}

Para que el usuario \texttt{tomasdea} pueda ejecutar comandos de Docker
sin \texttt{sudo} (requerido para que el servicio funcione sin elevar
privilegios), debe pertenecer al grupo \texttt{docker}:

\begin{verbatim}
sudo usermod -aG docker tomasdea
\end{verbatim}

El cambio de grupo toma efecto en la próxima sesión de login o reinicio.

% ─────────────────────────────────────────────────────────────────────────────
\subsection{Resumen de diferencias PC vs.\ Raspberry Pi}
\label{sec:docker-diff-summary}
% ─────────────────────────────────────────────────────────────────────────────

\begin{table}[H]
\centering
\begin{tabular}{lll}
\toprule
\textbf{Aspecto} & \textbf{PC (amd64)} & \textbf{Raspberry Pi (arm64)} \\
\midrule
Imagen Docker         & \texttt{tesis-ros2:humble-cpu}       & \texttt{tesis-ros2:humble-cpu-arm64} \\
Compose file          & \texttt{docker-compose.yml}           & \texttt{docker-compose.raspi.yml}    \\
librealsense          & Paquete apt v2.56                     & Compilada desde fuente v2.57.6       \\
Backend RealSense     & Kernel (V4L2/UVC)                     & RSUSB (espacio de usuario)           \\
Herramientas gráficas & RViz2, rqt, imshow                   & No instaladas                        \\
Display               & \texttt{/tmp/.X11-unix} montado       & No aplica                            \\
GPIO                  & No montado                            & \texttt{/dev/gpiochip0} + grupo 986  \\
\texttt{pid: host}    & No                                    & Sí (necesario para GPIO)             \\
\texttt{RMW}          & No fijado explícitamente              & \texttt{rmw\_cyclonedds\_cpp}        \\
Arranque automático   & Manual (\texttt{tesis-run})           & Servicio systemd                     \\
Build librealsense    & No aplica                             & \texttt{-j2} (límite de RAM)         \\
\bottomrule
\end{tabular}
\caption{Diferencias de configuración entre la PC de desarrollo y la
  Raspberry Pi de despliegue.}
\label{tab:docker-pc-vs-raspi}
\end{table}
